\subsection{代数可逆元群的李代数}

设 A 为带单位元 \(e\) 的完备可赋范结合代数。设 \(\mathbf{A}^*\) 为 A 的可逆元群。我们已看到 (§ 1, no. 1)，\(\mathbf{A}^*\) 是 A 的开子流形且是一个李群。设 \(\mathbf{G}\) 为李群，\(f\) 为李群 \(\mathbf{G}\) 到李群 \(\mathbf{A}^*\) 的态射，\(f\) 可以看作 \(\mathbf{G}\) 到完备可赋范空间 A 的解析映射。因而，若 \(t \in \mathcal{T}^{(\infty)}(\mathbf{G})\)，我们可以构造 \(\langle t, f \rangle\)，它是 A 的一个元素。

命题 33. 映射 \(t \mapsto \langle t, f \rangle\) 是代数 \(\mathcal{T}^{(\infty)}(G)\) 到代数 A 中的态射。

只需验证：若 \(t\) 与 \(t'\) 是 \(G\) 上的点分布，则 \(\langle t * t', f \rangle = \langle t, f \rangle \langle t', f \rangle\)。但

\[
\begin{array}{r l} \langle t * t ^ {\prime}, f \rangle & = \langle t \otimes t ^ {\prime}, (g, g ^ {\prime}) \mapsto f (g g ^ {\prime}) \rangle \\ & = \langle t \otimes t ^ {\prime}, (g, g ^ {\prime}) \mapsto f (g) f (g ^ {\prime}) \rangle \\ & = \langle t, f \rangle \langle t ^ {\prime}, f \rangle \\ & \quad (Diff. \& Anal. Man., R, 13.4.3). \end{array}
\]

命题 33 的态射称为与 f 相伴的态射。

取 G 为群 A* 自身，取 f 为 A* 的恒等映射 \(\iota\)。我们得到代数 \(\mathcal{T}^{(\infty)}(\mathbf{A}^*)\) 到代数 A 中的一个称为典范的态射。切空间 \(\mathrm{T}_{e}(\mathbf{A}^*)\) 典范地与 A 认同；且若 \(t \in \mathrm{T}_{e}(\mathbf{A}^*)\)，则这个认同的定义使得 \(\langle t, \iota \rangle = t\)。于是命题 33 蕴含下述推论：

推论。典范映射 \(\zeta\)（从 \(\mathbf{L}(\mathbf{A}^{*})\) 到 \(\mathbf{A}\)）是李代数 \(\mathbf{L}(\mathbf{A}^{*})\) 到李代数 \(\mathbf{A}\) 上的同构。换句话说，

\[
\zeta ([ a, b ]) = \zeta (a) \zeta (b) - \zeta (b) \zeta (a)
\]

对所有 \(a, b\)（\(\mathrm{L}(\mathbf{A}^*)\) 中的元素）成立。若 \(\mathbf{K}\) 的特征为 \(p > 0\)，则 \(\zeta(a^p) = \zeta(a)^p\)（对所有 \(a \in \mathrm{L}(\mathbf{A}^*)\)）。

今后我们借助同构 ζ 把 L(A*) 与 A 认同。

\(\mathcal{T}^{(\infty)}(\mathbf{A}^*)\) 到 \(\mathbf{A}\) 的典范态射是作为命题 33 的态射的特殊情形而得到的。但也可以沿相反的方向论证：

命题 34. 设 H 为李群，A 为含单位元的完备可赋范结合代数，\(\phi: \mathrm{H} \to \mathrm{A}^*\) 为李群态射。相伴态射 \(\phi'\)（\(\mathcal{T}^{(\infty)}(\mathrm{H})\) 到 A 中的）是由 \(\phi_*\) 与 \(\mathcal{T}^{(\infty)}(\mathrm{A}^*)\) 到 A 的典范态射复合而得。特别地，\(\phi'(x) = \mathrm{L}(\phi)(x)\)（对所有 \(x \in \mathrm{L}(\mathrm{H})\)）。

设 \(i\) 为 \(\mathbf{A}^*\) 到 A 的恒等映射。于是，对所有 \(t\in \mathcal{T}^{(\infty)}(\mathrm{H})\)，

\[
\begin{array}{r l} \phi^ {\prime} (t) = \langle t, \phi \rangle & = \langle t, i \circ \phi \rangle \\ & = \langle \phi_ {*} (t), i \rangle \quad (Diff. \& Anal. Man., R, 13.2.3). \end{array}
\]

